% Source-level obstruction to a proposed arithmetic completion.
% GPT-6 Astra Pro; reasoning effort not exposed. 2026-11-21.
\subsection{SF-Q6-07: the constant-twist row model cannot realize the forced term}

\paragraph{Statement SF-Q6-07.}
Suppose $D$ is completed from actual columns by
\texttt{RowsRestModel.model}, as in
\texttt{TrueColsTest.later}.  Let $w$ be a coefficient place above
a prime $p$ and let $1\le j\le m$.  Assume
\[
 p\notin\{g_0,g_1\},\qquad
 \texttt{IsResidueClassZero}(D,p),\qquad h_P(D,p)=0,
\]
and
\[
 k>1,\qquad \sum_{\tau}e_{j,\tau}>0.
\]
There is no $L:\texttt{SizeData}_2(D)$ simultaneously satisfying
$p\le X''$ and \texttt{L.FDef}.

\begin{proof}
The definition \texttt{RowsRestModel.model} sets
\texttt{muAt w := C.mu0} at every place.  The assembled record
retains this value.  Therefore the \texttt{untwisted} field of
any $L:\texttt{SizeData}_2(D)$ gives $F_j(w)=0$.

On the other hand \texttt{L.FDef}, the stated prime class, and
$p\le X''$ select the residue-class-zero branch of
\texttt{FvalP} and \texttt{FInert}, giving
\[
 F_j(w)=\frac{(k-1)\sum_\tau e_{j,\tau}}{p^2-1}>0.
\]
The denominator is positive because $p$ is prime.  This is a
contradiction.
\end{proof}

\paragraph{Application and scope.}
The proposed large-weight specialization in SF-Q6-04 has positive
$e_1=M+18$ and $e_2=M+54$ and is intended to use $p=863<X''$.
Thus the generic row-model completion cannot supply its arithmetic
record with the required forced-term definition.  This remains true
even though the column theta maps and their compact class-sum
periods inside that completion are actual maps: the auxiliary
row model supplies the wrong place-dependent twist support.

The obstruction does not apply to a different, actual row stage
whose \texttt{MuIs.on} field realizes the chosen twists on a finite
prime set containing $863$.  The datum
\texttt{TwDat.untw B pl} itself exists for every finite set
\texttt{pl}; constructing that datum is only one part of the
required row-stage and period-valuation realization.

\paragraph{Exact source locators.}
All paths are relative to \texttt{lean/src/}.
\begin{itemize}
\item \texttt{Development/AbelianVarietyTheory/Definitions/AutomorphicMembership/
ModelLamPDef.lean}, \texttt{RowsRestModel.model} (lines 130--135);
the actual compact-period formula is retained in its
\texttt{perCol} field (lines 162--166).
\item The same directory, \texttt{ExistsOkPins.lean},
\texttt{TrueColsTest.later} (lines 194--203), which selects
\texttt{RowsRestModel.model (fits l) 1 1 1};
\texttt{RowsRestWordColumnUnitBaseTuple.lean}, the assembly
retaining \texttt{muAt} (line 125).
\item \texttt{Formalism/LocalRingCardinality/DataParams$_2$.lean},
\texttt{SizeData$_2$.untwisted} (line 108).
\item \texttt{Formalism/LocalRingCardinality/DataForced.lean},
\texttt{FInert}, \texttt{FvalP}, \texttt{FDef}
(lines 112--135).
\item \texttt{Development/AbelianVarietyTheory/Definitions/AutomorphicMembership/
CmQuadraticFormEqAlternatingForm.lean}, \texttt{MuIs}
(lines 140--148) and \texttt{TwDat.untw} (lines 196--206).
\end{itemize}
The author read the complete defining model and the complete
\texttt{later} proof, and checked the assembly and forced-term
branches.  This is a human source-level argument; no Lean build is
claimed.
